\section{Results}
\paragraph{Release rule.}
Only values generated from frozen artifacts enter the final manuscript. Red dashes in this working revision identify required experiments; they are not zero values and must not be estimated manually.

\subsection{Data integrity and detector evidence}
The final version will first establish that source partitions are disjoint, that negative sources were reviewed, and that detector evidence is characterized by crop, object scale, and annotation density. Table~\ref{tab:data} will include the completed human-audit and detector-recall entries. Hard-negative mining will remain a supporting detector refinement unless its matched artifacts establish a reproducible gain.

\draftnote{Insert duplicate/leakage audit, human-review agreement, detector M$\rightarrow$M, A$\rightarrow$A, M$\rightarrow$A, and A$\rightarrow$M results.}

\subsection{Equal-budget evidence allocation}
Table~\ref{tab:allocation} is designed to report the complete policy trade-off rather than select a winner from one metric. The final analysis will compare paired source-level differences against combined top-$K$ and coordinate-aware spatial selection. A policy will be described as preferable only for an explicitly stated operating objective.

\begin{table*}[t]
\centering
\wacvtablefont
\caption{Equal-budget source-level comparison at the validation-selected primary configuration. Exact is strict source-record accuracy; AURC is lower-is-better. All intervals and paired comparisons are clustered by source.}
\label{tab:allocation}
\setlength{\tabcolsep}{5pt}
\begin{tabular}{lccccccc}
\toprule
Policy & Exact$\uparrow$ & Evid. rec.$\uparrow$ & Spatial cov.$\uparrow$ & Redund.$\downarrow$ & Hard-neg.$\uparrow$ & AURC$\downarrow$ & p95 ms$\downarrow$ \\
\midrule
Random & \pendingcell & \pendingcell & \pendingcell & \pendingcell & \pendingcell & \pendingcell & \pendingcell \\
Spatial (2D) & \pendingcell & \pendingcell & \pendingcell & \pendingcell & \pendingcell & \pendingcell & \pendingcell \\
Confidence top-$K$ & \pendingcell & \pendingcell & \pendingcell & \pendingcell & \pendingcell & \pendingcell & \pendingcell \\
Count top-$K$ & \pendingcell & \pendingcell & \pendingcell & \pendingcell & \pendingcell & \pendingcell & \pendingcell \\
Combined top-$K$ & \pendingcell & \pendingcell & \pendingcell & \pendingcell & \pendingcell & \pendingcell & \pendingcell \\
Exploit--explore & \pendingcell & \pendingcell & \pendingcell & \pendingcell & \pendingcell & \pendingcell & \pendingcell \\
Diversity-aware & \pendingcell & \pendingcell & \pendingcell & \pendingcell & \pendingcell & \pendingcell & \pendingcell \\
\bottomrule
\end{tabular}
\end{table*}

\begin{figure}[t]
\centering
\fbox{\parbox[c][36mm][c]{0.93\columnwidth}{\centering
\textbf{Pending artifact-generated policy Pareto plot}\\[2mm]
Source utility versus accepted risk; marker annotations report budget and latency.}}
\caption{\textbf{Allocation Pareto frontier.} The final figure will show paired source-level utility, risk, and compute for equal-budget routing policies.}
\label{fig:pareto}
\end{figure}

\subsection{Does the VLM add source-level value?}
Table~\ref{tab:vlmvalue} directly addresses the necessity of multimodal inference. Detector-only and deterministic aggregation are primary controls, not secondary ablations. The VLM claim will be retained only if source-level utility or selective behavior improves at a reported cost.

\begin{table*}[t]
\centering
\wacvtablefont
\caption{Source-level system comparison on the frozen pooled test protocol. Count error is reported continuously; FAR is false acceptance on verified negative sources.}
\label{tab:vlmvalue}
\setlength{\tabcolsep}{4.5pt}
\begin{tabular}{lcccccccc}
\toprule
System & Pres. F1 & Count err.$\downarrow$ & Dens. F1 & Cov. F1 & Exact$\uparrow$ & FAR$\downarrow$ & Review$\downarrow$ & p95 ms$\downarrow$ \\
\midrule
Detector only & \pendingcell & \pendingcell & \pendingcell & \pendingcell & \pendingcell & \pendingcell & -- & \pendingcell \\
Detector + rules & \pendingcell & \pendingcell & \pendingcell & \pendingcell & \pendingcell & \pendingcell & -- & \pendingcell \\
Whole-image VLM & \pendingcell & \pendingcell & \pendingcell & \pendingcell & \pendingcell & \pendingcell & \pendingcell & \pendingcell \\
Uniform tiled VLM & \pendingcell & \pendingcell & \pendingcell & \pendingcell & \pendingcell & \pendingcell & \pendingcell & \pendingcell \\
OrnaScout, no gate & \pendingcell & \pendingcell & \pendingcell & \pendingcell & \pendingcell & \pendingcell & -- & \pendingcell \\
OrnaScout + gate & \pendingcell & \pendingcell & \pendingcell & \pendingcell & \pendingcell & \pendingcell & \pendingcell & \pendingcell \\
\bottomrule
\end{tabular}
\end{table*}

\subsection{Directed crop transfer}
Table~\ref{tab:transfer} reports each training--test direction separately. The final analysis will give the paired drop from the corresponding in-domain reference and decompose it into detector miss, routing miss, structured-prediction error, and rejection behavior. Object scale and density are diagnostic strata, not post-hoc causal proof.

\begin{table}[t]
\centering
\wacvtablefont
\caption{Directed source-level transfer. Each cell reports the primary utility with source-clustered 95\% interval; $\Delta$ is the paired drop from the corresponding in-domain reference.}
\label{tab:transfer}
\setlength{\tabcolsep}{5pt}
\begin{tabular}{lcccc}
\toprule
Train & Test M & $\Delta$ & Test A & $\Delta$ \\
\midrule
Magnolia & \pendingcell & -- & \pendingcell & \pendingcell \\
Azalea & \pendingcell & \pendingcell & \pendingcell & -- \\
Pooled & \pendingcell & \pendingcell & \pendingcell & \pendingcell \\
\bottomrule
\end{tabular}
\end{table}

\subsection{Selective reliability and review burden}
Figure~\ref{fig:riskcoverage} will compare raw VLM confidence, strict-schema rejection, detector agreement, fixed-score ranking, and the cross-fitted learned gate. The operating point is selected from out-of-fold validation predictions before test scoring. The corresponding table reports coverage, accepted risk, negative-source false acceptance, accepted source count, and human-review rate.

\begin{figure}[t]
\centering
\fbox{\parbox[c][39mm][c]{0.93\columnwidth}{\centering
\textbf{Pending artifact-generated risk--coverage curves}\\[2mm]
Raw confidence, schema-only, detector agreement, fixed score, learned gate; validation-selected operating point marked.}}
\caption{\textbf{Held-out selective reliability.} Curves and intervals use source-level outcomes; lower risk at negligible coverage is not treated as operational success.}
\label{fig:riskcoverage}
\end{figure}

\begin{table}[t]
\centering
\wacvtablefont
\caption{Held-out reliability at the predeclared target-risk operating point.}
\label{tab:reliability}
\setlength{\tabcolsep}{3.4pt}
\begin{tabular}{lccccc}
\toprule
Score & AURC$\downarrow$ & Cov.$\uparrow$ & Risk$\downarrow$ & FAR$\downarrow$ & Accepted \\
\midrule
Raw confidence & \pendingcell & \pendingcell & \pendingcell & \pendingcell & \pendingcell \\
Strict schema & \pendingcell & \pendingcell & \pendingcell & \pendingcell & \pendingcell \\
Detector agreement & \pendingcell & \pendingcell & \pendingcell & \pendingcell & \pendingcell \\
Fixed score & \pendingcell & \pendingcell & \pendingcell & \pendingcell & \pendingcell \\
Learned gate & \pendingcell & \pendingcell & \pendingcell & \pendingcell & \pendingcell \\
\bottomrule
\end{tabular}
\end{table}

\subsection{Efficiency and failure analysis}
End-to-end timing will separate tiling, detection, selection, image transfer, VLM generation, parsing, gating, and reduction. The failure audit uses a frozen, crop-balanced, size-balanced source sample. Categories are annotation issue, detector miss, routing miss, boundary truncation, malformed schema, task error, unsupported positive, reduction conflict, gate false accept, gate false reject, and correct abstention.

\draftnote{Insert named-hardware efficiency table, failure counts with denominators and agreement, and frozen-rule qualitative panels.}
